\section{Conclusion}
\label{sec:conclusion}

We investigated how multilingual LLMs adjudicate direct factual contradictions between English and Korean retrieved evidence in mRAG. Introducing \textsc{MKQA-Conflict-150} and a counterbalanced $2 \times 2 \times 2$ protocol across query languages and passage orders isolated language bias across $4{,}200$ GPU runs: (1) \textbf{English Authority Bias ($\mathcal{H}_{\text{EN}}$):} Both models systematically favor English over Korean context ($\text{EPR}{=}\qwenEPRAll\%$ Qwen, $\chi^2{=}\qwenMcNemarChiAll$; $\llamaEPRAll\%$ Llama, $\chi^2{=}\llamaMcNemarChiAll, p \qwenMcNemarPAll$), with $>70\%$ of discordant pairs flipping toward English. (2) \textbf{Query-Language Moderation ($\mathcal{H}_{\text{QL}}$):} Korean queries reduce marginal $\text{EPR}$ in Qwen (\qwenEPREN\%$\to$\qwenEPRKO\%) and lower discordant flip skew in both models, yet $\text{EPR}_{\text{KO}}$ remains above $60\%$ ($p < 0.0001$), causing high Truth Override Rates ($\text{TOR}{=}\qwenPairTwoOverrideKO\%$ and $\llamaPairTwoOverrideKO\%$) under false English context. (3) \textbf{Partial Prior Protection ($\mathcal{H}_{\text{Prior}}$):} Parametric memory moderates truth override but fails to eliminate bias: false English context overrides truth in \qwenPKKnownBothOverrideAll\% (Qwen) and \llamaPKKnownBothOverrideAll\% (Llama) of $\text{Pair}_2$ runs on known facts. (4) \textbf{Script Anchoring vs.\ Authority:} English bias peaks on $\textsc{Date}$ (\qwenCatDateEPRAll\%/\llamaCatDateEPRAll\%) and $\textsc{Number}$ (\qwenCatNumberEPRAll\%/\llamaCatNumberEPRAll\%) where shared Arabic numerals offer no script cue, whereas Hangul output alignment partially anchors Korean queries on $\textsc{Entity}$ (\qwenCatEntityEPRAll\%/\llamaCatEntityEPRAll\%). Language fairness in mRAG thus requires generator-level interventions beyond retrieval and reranking \citep{park2025,wang2026}.
